%! Linear Regression — Unit 1, Lesson 1.5 — Warm-Up
\documentclass[12pt]{article}
\usepackage{algebra2-article}
\usepackage{algebra2-boxes}
% Coordinate grid: {xmin}{xmax}{ymin}{ymax}
\newcommand{\lgrid}[4]{%
  \draw[step=1, linegray!40, very thin] (#1,#3) grid (#2,#4);
  \draw[->, semithick, charcoal] (#1,0)--(#2,0) node[below right, font=\scriptsize]{$x$};
  \draw[->, semithick, charcoal] (0,#3)--(0,#4) node[left, font=\scriptsize]{$y$};}
\newcommand{\dpt}[2]{\fill[navy] (#1,#2) circle (3pt);}

\begin{document}

\pageheader{Unit 1, Lesson 1.5}{Warm-Up}
% NAMESTRIP: no \namedateperiod here — the student writes their name once, on the cover.

\vspace{0.10in}

\begin{notesbox}{Getting Ready: Skills We'll Use Today}
Three things you already know. Today we work with \emph{scattered} real data; these are the
tools you'll need --- work quickly, then check with a neighbour.

\vspace{6pt}
\textbf{1. Read the points.} Give the coordinates $(x,y)$ of each labelled point.
\begin{center}
\begin{tikzpicture}[scale=0.46]
  \lgrid{-0.5}{7}{-0.5}{6}
  \dpt{1}{2} \node[font=\scriptsize\bfseries, navy] at (1,2.8){A};
  \dpt{3}{3} \node[font=\scriptsize\bfseries, navy] at (3,3.8){B};
  \dpt{5}{5} \node[font=\scriptsize\bfseries, navy] at (5,5.8){C};
  \dpt{6}{2} \node[font=\scriptsize\bfseries, navy] at (6,2.8){D};
\end{tikzpicture}
\end{center}
\vspace{-4pt}
$A=$ \blank{1.4cm} \quad $B=$ \blank{1.4cm} \quad $C=$ \blank{1.4cm} \quad $D=$ \blank{1.4cm}
\\[4pt]
Do the four points lie on one straight line? \blank{1.2cm} \; As $x$ gets bigger, do the
$y$-values \emph{mostly} rise, fall, or neither? \blank{1.8cm}

\vspace{8pt}
\textbf{2. Predict with a rule.} A line is $y=4x+10$. Evaluate it at each input.
\begin{center}
\renewcommand{\arraystretch}{1.3}
\begin{tabular}{|c|c|c|c|}
\hline
$x$ & $2$ & $5$ & $8$ \\ \hline
$y=4x+10$ & \blank{1.3cm} & \blank{1.3cm} & \blank{1.3cm} \\ \hline
\end{tabular}
\end{center}

\vspace{6pt}
\textbf{3. Slope and direction.} For $y=ax+b$, the slope is $a$ and the $y$-intercept is $b$.
\begin{itemize}[leftmargin=1.6em, itemsep=4pt]
  \item For $y=3x-5$: \; slope $=$ \blank{1.2cm}, \; $y$-intercept $=$ \blank{1.2cm}.
  \item A line with a \emph{positive} slope goes \blank{2.0cm} (up / down) as you read left to
        right; a \emph{negative} slope goes \blank{2.0cm}.
  \item In $y=3x-5$, each time $x$ goes up by $1$, $y$ goes up by \blank{1.2cm}.
\end{itemize}
\end{notesbox}

\end{document}
